\documentclass[11pt,a4paper]{article}
\usepackage[utf8]{inputenc}
\usepackage[T1]{fontenc}
\usepackage[french]{babel}
\usepackage[top=2cm,bottom=2cm,left=2cm,right=2cm]{geometry}
\usepackage{xcolor}
\usepackage{tcolorbox}
\tcbuselibrary{skins,breakable}
\usepackage{fancyhdr}
\usepackage{enumitem}
\usepackage{lmodern}

% --- Couleur discipline : MANAGEMENT ---
\definecolor{mainColor}{RGB}{0,120,60}
\definecolor{darkGray}{RGB}{50,50,50}

\tcbset{boxrule=0.5pt, arc=3pt, left=5pt, right=5pt, top=3pt, bottom=3pt, breakable}

\newtcolorbox{defbox}[1]{
  colback=mainColor!8, colframe=mainColor, coltitle=white,
  fonttitle=\bfseries\footnotesize, title=DÉFINITION \textemdash\ #1}

\newtcolorbox{keybox}{
  colback=darkGray!7, colframe=darkGray!50, coltitle=white,
  fonttitle=\bfseries\footnotesize, title=POINT CLÉ}

\newtcolorbox{exbox}{
  colback=mainColor!5, colframe=mainColor!40,
  fonttitle=\bfseries\footnotesize, title=EXEMPLE}

\pagestyle{fancy}\fancyhf{}
\fancyhead[L]{\small\textcolor{mainColor}{\textbf{CEJM — BTS SIO}}}
\fancyhead[C]{\small\textbf{Chapitre 9 — L'environnement global de l'entreprise}}
\fancyhead[R]{\small\textcolor{mainColor}{\textbf{MANAGEMENT}}}
\fancyfoot[C]{\small\thepage}
\renewcommand{\headrulewidth}{0.5pt}
\setlength{\headheight}{15pt}
\setlist{itemsep=2pt, topsep=3pt, parsep=0pt}

\begin{document}

\begin{center}
  {\Large\bfseries\textcolor{mainColor}{Chapitre 9}}\\[2pt]
  {\large L'environnement global de l'entreprise}\\[4pt]
  \textit{\textcolor{darkGray}{Comment l'entreprise analyse-t-elle son environnement pour prendre ses décisions ?}}
\end{center}
\vspace{4pt}\hrule\vspace{8pt}

\section*{Définitions essentielles}

\begin{defbox}{Macro-environnement}
\textbf{Environnement général} de l'entreprise : facteurs sur lesquels elle a peu ou pas d'influence (tendances sociétales, réglementations, conjoncture).
\end{defbox}

\vspace{4pt}
\begin{defbox}{Micro-environnement}
\textbf{Environnement proche} de l'entreprise, avec lequel elle interagit directement : clients, concurrents, fournisseurs, partenaires.
\end{defbox}

\vspace{4pt}
\begin{defbox}{Analyse PESTEL}
Outil de diagnostic de l'\textbf{environnement externe} qui classe les facteurs en 6 axes : Politique, Économique, Socioculturel, Technologique, Écologique, Légal.
\end{defbox}

\vspace{4pt}
\begin{defbox}{Innovation}
Introduction d'une \textbf{nouveauté} (produit, procédé, organisation, marché) permettant à l'entreprise de s'adapter aux évolutions de l'environnement et de se différencier.
\end{defbox}

\section*{Macro vs Micro-environnement}

\begin{center}
\begin{tabular}{|l|l|l|}
\hline
 & \textbf{Macro-environnement} & \textbf{Micro-environnement} \\
\hline
\textbf{Influence de l'entreprise} & Très faible & Directe (possible) \\
\textbf{Exemples} & Lois, conjoncture, climat & Clients, concurrents, fournisseurs \\
\textbf{Outils d'analyse} & PESTEL & 5 forces de Porter \\
\hline
\end{tabular}
\end{center}

\section*{L'analyse PESTEL}

\begin{keybox}
Chaque facteur PESTEL peut être une \textbf{opportunité à saisir} ou une \textbf{menace à anticiper}.
\end{keybox}

\vspace{4pt}
\begin{tabular}{|c|l|l|}
\hline
\textbf{Lettre} & \textbf{Facteur} & \textbf{Exemples} \\
\hline
\textbf{P} & Politique & Stabilité gouvernementale, politique fiscale, commerce extérieur \\
\textbf{E} & Économique & Croissance du PIB, inflation, taux d'intérêt, chômage, pouvoir d'achat \\
\textbf{S} & Socioculturel & Vieillissement pop., modes de vie, niveau d'éducation, consumérisme \\
\textbf{T} & Technologique & R\&D, digitalisation, brevets, IA, infrastructures numériques \\
\textbf{E} & Écologique & Changement climatique, normes environnementales, énergies renouvelables \\
\textbf{L} & Légal & Droit du travail, droit de la concurrence, RGPD, normes produits \\
\hline
\end{tabular}

\vspace{6pt}
\begin{exbox}
\textbf{Entreprise de e-commerce} : Opportunité \textbf{T} (essor du mobile) + Opportunité \textbf{E} (pouvoir d'achat stable). Menace \textbf{L} (RGPD contraignant) + Menace \textbf{E} (inflation = baisse des achats).
\end{exbox}

\section*{Les conséquences sur la prise de décision}

\begin{itemize}
  \item \textbf{Opportunités} : évolutions favorables à saisir (nouveau marché, nouvelle techno, déréglementation)
  \item \textbf{Menaces} : contraintes à anticiper (crise économique, réglementation stricte, rupture technologique)
  \item \textbf{Incertitude} : l'environnement est complexe et volatile → la décision se prend sous contrainte
\end{itemize}

\section*{L'innovation comme réponse aux évolutions}

\begin{keybox}
Innover permet de \textbf{transformer une menace en opportunité} et d'assurer la pérennité de l'entreprise.
\end{keybox}

\begin{center}
\begin{tabular}{|l|l|}
\hline
\textbf{Type d'innovation} & \textbf{Description} \\
\hline
Produit & Nouveau bien ou service (ex. : smartphone) \\
Procédé & Nouvelle méthode de production (ex. : impression 3D) \\
Organisationnelle & Nouvelle façon d'organiser le travail (ex. : télétravail) \\
Commerciale & Nouveau canal de vente (ex. : marketplace) \\
\hline
\end{tabular}
\end{center}

\begin{exbox}
Le \textbf{changement climatique} (facteur E écologique) est une menace pour l'industrie automobile thermique → opportunité pour les constructeurs de véhicules électriques.
\end{exbox}

\section*{Points clés à retenir}

\begin{keybox}
L'\textbf{analyse PESTEL} est l'outil de référence pour scanner le macro-environnement.
\end{keybox}
\vspace{3pt}
\begin{keybox}
L'environnement génère des \textbf{opportunités et des menaces} qu'il faut identifier pour adapter la stratégie.
\end{keybox}
\vspace{3pt}
\begin{keybox}
L'\textbf{innovation} est le principal levier pour s'adapter aux évolutions et rester compétitif.
\end{keybox}

\section*{Mots-clés}
\textcolor{mainColor}{\textbf{Macro-environnement}} $\cdot$
\textcolor{mainColor}{\textbf{Micro-environnement}} $\cdot$
\textcolor{mainColor}{\textbf{PESTEL}} $\cdot$
\textcolor{mainColor}{\textbf{Opportunité / Menace}} $\cdot$
\textcolor{mainColor}{\textbf{Innovation}} $\cdot$
\textcolor{mainColor}{\textbf{Prise de décision}} $\cdot$
\textcolor{mainColor}{\textbf{Veille stratégique}} $\cdot$
\textcolor{mainColor}{\textbf{Mondialisation}} $\cdot$
\textcolor{mainColor}{\textbf{Digitalisation}}

\end{document}
